\section{Cái sổ, và chữ giao}
\label{sec:ch18-so-va-chu-giao}

\index{khoảng trống nghiên cứu!cái sổ gồm những gì|(}%
Phần mở đầu chương vừa nói chữ \emph{giao} sai ngay từ kế hoạch. Mục này nói cái
sổ thật ra chứa gì, vì phải biết tập hợp trước rồi mới hỏi được nó có giao hay
không.

\index{khoảng trống nghiên cứu!chín bài của Phần IV}%
Cái sổ được lập theo một quy tắc đặt ra từ chương~\ref{ch:do-do-trung-thuc}: mỗi
phiên viết một chương của Phần~IV ghi lại toàn bộ phát biểu hạn chế và toàn bộ
phần công việc tương lai của bài mình đọc. Bảy chương của Phần~IV đọc chín bài
báo, vì chương~\ref{ch:con-nguoi-vang-mat} đọc hai bài và
chương~\ref{ch:giai-thich-duoi-tan-cong} đọc ba;
chương~\ref{ch:chi-so-khong-do-do-trung-thuc} không thêm mục nào vào sổ vì bài nó
đọc đã được chương~\ref{ch:do-do-trung-thuc} ghi. Chín bài ấy là toàn bộ cái sổ
mà chương này thừa kế.

Ba chương của Phần~V ghi thêm sáu bài nữa, và ghi vào chỗ khác chứ không vào sổ
ấy. Đó là chủ ý: cái sổ được lập cho các bài của Phần~IV, còn các bài của Phần~V được
đọc để trả lời câu hỏi lối thoát chứ không để đếm chỗ thiếu. Chương này đọc cả
mười lăm bài, và giữ hai nhóm tách nhau ở chỗ nào có ích.

\index{khoảng trống nghiên cứu!sáu bài giao và ba bài không}%
Bây giờ hỏi cái sổ có giao hay không. Sáu trong chín bài nêu cùng một thứ đang
thiếu, là một dụng cụ đo \tn{độ trung thực}{faithfulness} đã được kiểm định, và
nêu từ sáu hướng
khác nhau: đo thấy các dụng cụ hỏng, ghi nhận rằng không có chuẩn nào để dùng
chúng, gọi tên bước tiếp theo mà chính tác giả chưa làm, đếm xem phép đối chiếu
với người hiếm tới mức nào, chạy phép đối chiếu ấy một lần rồi thấy nó không
khớp, và tự nhận chưa đo xem công cụ mình vừa đề xuất có giải thích được gì
không. Ba bài còn lại không nêu, và đây là chỗ mà chữ \emph{giao} hỏng: ba chỗ
trống của ba bài ấy không giống nhau và cũng không giống sáu bài kia. Bài của
chương~\ref{ch:gioi-han-ly-thuyet} muốn một đại lượng tính được thay cho một đại
lượng không tính được. Bài của mục~\ref{sec:ch14-shap-lam-vu-khi} muốn những mô
hình khái quát hóa tốt hơn. Bài của mục~\ref{sec:ch14-sua-bang-nhan-qua} muốn
thuật toán \tn{khám phá nhân quả}{causal discovery} tốt hơn.

Lấy giao của chín tập hợp mà ba trong số đó không chứa phần tử chung nào thì được
tập rỗng. Vậy nếu đọc chữ \emph{giao} theo nghĩa đen thì kế hoạch tự bác nó ngay
khi bài thứ nhất trong ba bài kia được đọc, và đó là bài của
chương~\ref{ch:gioi-han-ly-thuyet}, đọc từ giữa Phần~IV.

\index{khoảng trống nghiên cứu!ba chỗ không giao nói được gì}%
Ba chỗ không giao ấy vẫn là dữ kiện, và chúng nói một điều mà sáu bài kia không
nói được. Một bài không nêu chỗ thiếu vì hai lý do khác nhau: hoặc bài ấy không
đứng ở chỗ có thể thấy nó, hoặc bài ấy đứng ở đó mà nhìn sang hướng khác. Bài
của chương~\ref{ch:gioi-han-ly-thuyet} thuộc loại thứ nhất, và
mục~\ref{sec:ch13-tran-noi-gi} đã nói vì sao: đối tượng của nó là
độ khớp chứ không phải độ trung thực, nên chỗ thiếu của nó nằm ở
chuỗi khác. Hai bài kia thuộc loại thứ hai. Cả hai làm việc trên attribution
\tn{đặc trưng}{feature}, cả hai chấm điểm kết quả của mình bằng một con số, và
phần công việc
tương lai của cả hai đòi cải thiện đúng thứ mà bài ấy vốn đã dựng. Đó không phải
sơ suất của hai bài; đó là hình dạng mà mục sau đo được trên cả chín bài cùng
lúc.

\index{khoảng trống nghiên cứu!cái sổ gồm những gì|)}%
Vậy cái mà chương này lấy ra khỏi sổ không phải một giao. Nó là một hình, và hình
ấy có ba phần: sáu bài nói cùng một chỗ thiếu, ba bài nói chỗ khác, và một khoảng
cách giữa điều các bài chẩn đoán với điều các bài đề nghị. Phần thứ ba là phần mà
cái sổ chưa đo, vì cái sổ được lập bằng cách hỏi từng bài đang thiếu gì chứ không
hỏi nó định làm gì.
